% Options for packages loaded elsewhere
\PassOptionsToPackage{unicode}{hyperref}
\PassOptionsToPackage{hyphens}{url}
\documentclass[
  9pt,
]{article}
\usepackage{xcolor}
\usepackage[left=1.5cm,right=1.5cm,top=1.2cm,bottom=1.6cm]{geometry}
\usepackage{amsmath,amssymb}
\setcounter{secnumdepth}{-\maxdimen} % remove section numbering
\usepackage{iftex}
\ifPDFTeX
  \usepackage[T1]{fontenc}
  \usepackage[utf8]{inputenc}
  \usepackage{textcomp} % provide euro and other symbols
\else % if luatex or xetex
  \usepackage{unicode-math} % this also loads fontspec
  \defaultfontfeatures{Scale=MatchLowercase}
  \defaultfontfeatures[\rmfamily]{Ligatures=TeX,Scale=1}
\fi
\usepackage{lmodern}
\ifPDFTeX\else
  % xetex/luatex font selection
\fi
% Use upquote if available, for straight quotes in verbatim environments
\IfFileExists{upquote.sty}{\usepackage{upquote}}{}
\IfFileExists{microtype.sty}{% use microtype if available
  \usepackage[]{microtype}
  \UseMicrotypeSet[protrusion]{basicmath} % disable protrusion for tt fonts
}{}
\makeatletter
\@ifundefined{KOMAClassName}{% if non-KOMA class
  \IfFileExists{parskip.sty}{%
    \usepackage{parskip}
  }{% else
    \setlength{\parindent}{0pt}
    \setlength{\parskip}{6pt plus 2pt minus 1pt}}
}{% if KOMA class
  \KOMAoptions{parskip=half}}
\makeatother
\usepackage{graphicx}
\makeatletter
\newsavebox\pandoc@box
\newcommand*\pandocbounded[1]{% scales image to fit in text height/width
  \sbox\pandoc@box{#1}%
  \Gscale@div\@tempa{\textheight}{\dimexpr\ht\pandoc@box+\dp\pandoc@box\relax}%
  \Gscale@div\@tempb{\linewidth}{\wd\pandoc@box}%
  \ifdim\@tempb\p@<\@tempa\p@\let\@tempa\@tempb\fi% select the smaller of both
  \ifdim\@tempa\p@<\p@\scalebox{\@tempa}{\usebox\pandoc@box}%
  \else\usebox{\pandoc@box}%
  \fi%
}
% Set default figure placement to htbp
\def\fps@figure{htbp}
\makeatother
\setlength{\emergencystretch}{3em} % prevent overfull lines
\providecommand{\tightlist}{%
  \setlength{\itemsep}{0pt}\setlength{\parskip}{0pt}}
%% tools/stamp.tex
%% Provenance footer for PDFs rendered from this project.
%% Vendored by zzcollab; this file is not project-specific. The
%% footer prints the source document and the time of rendering.
%% The git version is defined here too, and so is available to a
%% project that wants it on the page, but it is left off the
%% default footer: it already names the staged copy in share/ and
%% appears in its manifest row. All values are supplied at render
%% time by tools/stamp-render.R, which writes a small generated
%% file that redefines the macros below. The \providecommand
%% defaults keep a bare LaTeX compile working when the document is
%% built without the wrapper.
\usepackage{fancyhdr}
\usepackage{xcolor}
\providecommand{\stampsource}{\detokenize{(source not recorded)}}
\providecommand{\stamptime}{(time not recorded)}
\providecommand{\stampversion}{\detokenize{(version not recorded)}}
\setlength{\footskip}{40pt}
\newcommand{\stampfootcontent}{%
  \footnotesize\color{gray}%
  \parbox[b]{\textwidth}{%
    \stamptime\hfill\thepage\\[2pt]%
    \ttfamily\raggedright\stampsource}}
\pagestyle{fancy}
\fancyhf{}
\fancyfoot[C]{\stampfootcontent}
\renewcommand{\headrulewidth}{0pt}
\renewcommand{\footrulewidth}{0.4pt}
%% The title page uses the `plain` page style; redefine it so the
%% stamp appears there too.
\fancypagestyle{plain}{%
  \fancyhf{}%
  \fancyfoot[C]{\stampfootcontent}%
  \renewcommand{\headrulewidth}{0pt}%
  \renewcommand{\footrulewidth}{0.4pt}}
\renewcommand{\stampsource}{\textasciitilde{}/\allowbreak{}prj/\allowbreak{}res/\allowbreak{}36-pmsimstats-ng/\allowbreak{}pmsimstats-ng/\allowbreak{}analysis/\allowbreak{}report/\allowbreak{}02-carryover-sensitivity/\allowbreak{}whitepaper-g9.Rmd}
\renewcommand{\stamptime}{2026-10-28 11:15 PDT}
\renewcommand{\stampversion}{\detokenize{whitepapers-pre-series-108-g10d5f1464-wip}}
\usepackage{amsmath}
\usepackage{booktabs}
\usepackage{titling}
\pretitle{\vspace{-1.2cm}\begin{center}\large\bfseries}
\posttitle{\end{center}\vspace{-0.9cm}}
\preauthor{\begin{center}\small}
\postauthor{\end{center}\vspace{-0.6cm}}
\predate{}
\postdate{}
\setlength{\parskip}{2pt}
\setlength{\parindent}{0pt}
\usepackage{booktabs}
\usepackage{longtable}
\usepackage{array}
\usepackage{multirow}
\usepackage{wrapfig}
\usepackage{float}
\usepackage{colortbl}
\usepackage{pdflscape}
\usepackage{tabu}
\usepackage{threeparttable}
\usepackage{threeparttablex}
\usepackage[normalem]{ulem}
\usepackage{makecell}
\usepackage{xcolor}
\usepackage{bookmark}
\IfFileExists{xurl.sty}{\usepackage{xurl}}{} % add URL line breaks if available
\urlstyle{same}
% fallback for those not using the hyperref driver hyperxmp:
\makeatletter
\@ifundefined{xmpquote}{\newcommand{\xmpquote}[1]{#1}}{}
\makeatother
\hypersetup{
  pdftitle={Analysis G9: testing a biomarker-by-drug interaction in an aggregated N-of-1 trial with carryover},
  pdfauthor={pmsimstats team},
  hidelinks,
  pdfcreator={LaTeX via pandoc}}

\title{Analysis G9: testing a biomarker-by-drug interaction in an
aggregated N-of-1 trial with carryover}
\author{pmsimstats team}
\date{}

\begin{document}
\maketitle

\textbf{Purpose.} G9 is the analysis we recommend for the
biomarker-by-drug interaction in a Hybrid aggregated N-of-1 trial when
carryover is present and the carryover half-life is unknown. It was the
best of the analyses evaluated in a staged simulation under the
correlation encoding of the interaction.

\textbf{Model.} For participant \(i\) at visit \(t\), with \(w_{it}\)
the week of the visit,
\[Y_{it}=\beta_0+\beta_B B_i+\beta_w w_{it}+\beta_D D_{bc,it}(h)+\beta_{BD}\,B_i\,D_{bc,it}(h)+u_i+\varepsilon_{it},\]
\[u_i\sim N(0,\sigma_u^2),\qquad \operatorname{Corr}(\varepsilon_{it},\varepsilon_{is})=\varphi^{|w_{it}-w_{is}|},\qquad
D_{bc,it}(h)=\begin{cases}1&\text{on drug}\\ 2^{-t_{sd,it}/h}&\text{off drug after exposure}\\ 0&\text{before first exposure.}\end{cases}\]
The half-life is chosen from the data,
\(\hat h=\arg\min_{h\in\mathcal T}\mathrm{AIC}\{M(h)\}\) with
\(\mathcal T=\{0.25,0.5,1,2\}\) weeks and AIC from maximum-likelihood
fits. The selected model is refitted by restricted maximum likelihood.
The interaction is tested with
\(T=\hat\beta_{BD}/\mathrm{SE}_{CR2}(\hat\beta_{BD})\) against a \(t\)
distribution with Satterthwaite degrees of freedom, where
\[\mathrm{SE}^2_{CR2}=\Bigl[(X^\top WX)^{-1}\sum_i X_i^\top W A_i\,\hat e_i\hat e_i^\top A_i^\top W X_i\,(X^\top WX)^{-1}\Bigr]_{BD,BD},\qquad A_i=(I-H_{ii})^{-1/2}.\]

\textbf{Notation.} \(B_i\): baseline biomarker, centered. \(D_{bc,it}\):
exposure-decayed drug indicator. \(t_{sd,it}\): weeks since the last
discontinuation (0 before exposure). \(h\): analysis half-life;
\(\hat h\): its selected value. \(u_i\): participant random intercept.
\(\varepsilon_{it}\): residual with continuous-time AR(1) correlation
\(\varphi\). \(\beta_{BD}\): the interaction, the change in drug effect
per unit of biomarker (under the correlation encoding
\(-c_{bm}\sigma_{BR}\) per standard deviation of the biomarker at the
true half-life). \(X_i\), \(W\), \(\hat e_i\), \(H_{ii}\): participant
design matrix, working weights, residuals, and leverage block. CR2: the
bias-reduced cluster-robust variance estimator (Bell and McCaffrey).

\textbf{Why G9.} The ladder adds one change at a time (original analysis
E0; G1 binary drug with AR(1); G6 adds CR2; G8 adds exposure decay at
the true half-life; G9 selects the half-life). Pooled over the corrected
original grid (Hybrid, OL+BDC and CO; \(N\) of 35 and 70):

\begin{table}
\centering\begingroup\fontsize{8}{10}\selectfont

\begin{tabular}{llll}
\toprule
Analysis & What changes & Null rejection (MCSE) & Power at $c_{bm}=0.45$, $t_{1/2}=1$ (MCSE)\\
\midrule
E0 & Original analysis & 0.110 (0.006) & 0.691 (0.008)\\
G1 & Binary drug, AR(1) & 0.038 (0.003) & 0.563 (0.009)\\
G6 & + CR2 inference & 0.055 (0.004) & 0.587 (0.009)\\
G8 & + exposure decay, true $h$ & 0.050 (0.004) & 0.640 (0.009)\\
G9 & + selected $\hat h$ (recommended) & 0.056 (0.004) & 0.641 (0.009)\\
F1 & Expectancy terms, fixed $h$ & 0.053 (0.005) & 0.517 (0.011)\\
\bottomrule
\end{tabular}
\endgroup{}
\end{table}

\textbf{Scenarios tested} (G9 in every one; each with a null and a power
cell unless noted). Designs: Hybrid, OL+BDC and crossover, \(N\) of 35
and 70, half-lives 0, 0.5 and 1 week. Misspecified decay: Weibull shapes
0.5 and 2 analyzed as exponential; true half-lives 0.2, 2 and 4 weeks.
Serial correlation 0.5 and 0.9. Dropout: completely at random and
dependent on baseline severity, 10, 20 and 30 percent. Prognostic
biomarkers: correlation with the placebo response, the natural course
and baseline severity, 0.3 and a feasible ladder to 0.8. Large \(N\) of
140 and 280 (null only). A biomarker that predicts the natural-history
trend and a patient-specific slope. A Hybrid run with eight null and
eight power cells at 5,000 and 1,000 replicates.

\begin{itemize}
\tightlist
\item
  \textbf{Size.} G9 rejects at 0.056 under the null, against 0.110 for
  E0, which is liberal, and 0.038 for G1, which is conservative. Across
  seven scenario families (misspecified decay, serial correlation,
  dropout, prognostic pathways, large \(N\), trend) its pooled size is
  0.045 to 0.060.
\item
  \textbf{Power.} G9 has 0.641 power, the same as G8 (0.640), which
  needs the true half-life. It exceeds G1 (0.563) and F1 (0.517). E0 has
  higher raw power only because it rejects too often.
\item
  \textbf{No pharmacokinetic input.} The analyst supplies a grid, not a
  half-life. Power matches the true-half-life analysis even though the
  selected value is spread over the grid.
\item
  \textbf{Costs, stated.} Selection adds 0.0059 (paired SE 0.0013) to
  size over G8 (0.055 against 0.049 pooled over the null cells), and its
  95\% interval covers about 0.93 to 0.94 of the time. A fixed 1-week
  half-life with CR2 avoids selection, holds nominal size, and stays
  within two points of G8's power.
\item
  \textbf{Limits.} The results hold under the correlation encoding, in
  the Hybrid design, on one parameter set. A biomarker that predicts the
  natural-history trend inflates the size of every analysis examined; G9
  rejects at 0.074 at \(N=70\).
\end{itemize}

\end{document}
